\documentclass[tikz]{standalone}
\usetikzlibrary{positioning, arrows.meta, calc, fit, backgrounds}
\usetikzlibrary{shapes.geometric}

\begin{document}
	
	\begin{tikzpicture}[
		% Ajustamos el espaciado interno por defecto a 1.8cm para los bloques hermanos
		node distance=2.0cm and 1.8cm,
		% Estilos base homogéneos
		bloque/.style={draw=black!50, fill=white, thick, rectangle, rounded corners, minimum height=1.8cm, minimum width=3.2cm, align=center, font=\sffamily},
		archivo/.style={draw=orange!70!black, fill=orange!10, thick, cylinder, shape border rotate=90, minimum height=1.2cm, minimum width=1.2cm, aspect=0.3, align=center, font=\scriptsize\sffamily},
		visualizador/.style={draw=black!50, fill=white, thick, rectangle, rounded corners, align=center, font=\sffamily, inner sep=0.15cm},
		env_box/.style={draw, thick, rounded corners=12pt, inner sep=15pt, fill opacity=0.04},
		flecha/.style={-{Stealth[scale=1.2]}, thick, draw=gray!80!black},
		etiqueta/.style={align=center, font=\tiny\sffamily, fill=white, inner sep=2.5pt, rounded corners=2pt}
		]
		
		% =================================================================
		% --- FILA SUPERIOR: TRACK DE PROCESAMIENTO PRINCIPAL -------------
		% =================================================================
		
		% C1_TOP: Simulink 
		\node[bloque] (matlab) {
			
\includegraphics[height=0.9cm, keepaspectratio]{Pictures/matlab_logo.png} \\[0.05cm]
			\textbf{Simulink}\\
			\scriptsize Motor de Física
		};
		
		% C2_TOP: Tópico ROS 2 (Aumentamos a 3.8cm para separar el bloque verde del azul)
		\node[bloque, right=3.8cm of matlab] (ros) {
			
\includegraphics[height=0.6cm, keepaspectratio]{Pictures/Ros_logo.png} \\[0.05cm]
			\scriptsize Topic:\\
			\texttt{/joint\_states}
		};
		
		% C3_TOP: Robot State Publisher
		\node[bloque, right=of ros, minimum width=2.8cm, dashed] (rsp) {
			robot\_state\_\\ publisher
		};
		
		% C4_TOP: Visualizador RViz
		\node[visualizador, right=of rsp] (rviz) {
			\includegraphics[width=3.2cm]{Pictures/modelo_3D} \\
			\footnotesize \textbf{RViz}
		};
		
		% =================================================================
		% --- FILA INFERIOR: REJILLA DE ALINEACIÓN ABSOLUTA ---------------
		% =================================================================
		
		% C1_BOT: Microcontrolador (Logo ampliado a 0.8cm para que no se vea enano)
		\node[bloque, below=of matlab, draw=red!70!black, fill=red!5] (uc) {
			\includegraphics[height=0.8cm, keepaspectratio]{Pictures/ti_logo.png} \\[0.1cm]
			\textbf{Microcontrolador}\\
			\scriptsize Planta Feedback 33-005
		};
		
		% C2_BOT: Nodo Serie (Se alinea automáticamente con 'ros' manteniendo la separación con 'uc')
		\node[bloque, draw=green!60!black] (serial_node) at (ros |- uc) {
			\textbf{Nodo Serie ROS 2}\\
			\scriptsize Pasarela UART $\rightarrow$ ROS\\
			\tiny (Python / C++)
		};
		
		% C3_BOT: Archivo URDF
		\node[archivo] (urdf) at (rsp |- uc) {Archivo\\URDF};
		
		% =================================================================
		% --- CAPAS DE FONDO (Contenedores) -------------------------------
		% =================================================================
		\begin{scope}[on background layer]
			% Caja Simulación Virtual (Azul)
			\node[env_box, draw=blue!70, fill=blue, fit=(matlab)] (caja_matlab) {};
			\node[above=0.2cm of caja_matlab.north, font=\sffamily\bfseries\color{blue!70!black}] {Entorno de Simulación Virtual};
			
			% Caja Hardware Real (Roja)
			\node[env_box, draw=red!70, fill=red, fit=(uc)] (caja_hw) {};
			\node[above=0.2cm of caja_hw.north, font=\sffamily\bfseries\color{red!70!black}] {Entorno de Hardware Real};
			
			% Caja Ecosistema ROS 2 (Verde)
			\node[env_box, draw=green!60!black, fill=green, fit=(serial_node) (ros) (rsp) (rviz) (urdf)] (caja_ros) {};
			\node[above=0.2cm of caja_ros.north, font=\sffamily\bfseries\color{green!60!black}] {Visualización del Gemelo Digital (Ecosistema ROS 2)};
		\end{scope}
		
		% =================================================================
		% --- CONEXIONES Y FLUJOS DE DATOS --------------------------------
		% =================================================================
		
		% Flujos Horizontales 
		\draw[flecha] (matlab.east) -- node[etiqueta, above=0.08cm] {ROS Toolbox} (ros.west);
		\draw[flecha] (uc.east) -- node[etiqueta, above=0.08cm] {SCI / UART} (serial_node.west);
		
		% Flujos Verticales
		\draw[flecha] (serial_node.north) -- node[etiqueta, xshift=0.15cm, pos=0.5, anchor=west] {} (ros.south);
		\draw[flecha, dashed, draw=orange!70!black] (urdf.north) -- (rsp.south);
		
		% Flujo hacia RViz
		\draw[flecha] (ros.east) -- (rsp.west);
		\draw[flecha] (rsp.east) -- node[etiqueta, above=0.08cm] {TF} (rviz.west);
		
	\end{tikzpicture}
	
\end{document}